\section{Open questions and next steps}
\label{sec:open}

\begin{itemize}
  \item \textbf{Other bases.} The construction is base-independent. Determine, for the bases covered
        by~\cite{dbts}, whether the infinite path of the child graph is unique and whether its
        branch choices are eventually periodic. The periods $L(b)$ grow quickly ($48\,720$ for
        $b = 11$), so this needs a faster implementation, for example one based on~\cite{dbtscode}.
  \item \textbf{A proof without the computer.} Explain the unique immortal entry per decade, and the
        pattern of choices, from the digit dynamics, as \cite[Theorem~10.1]{comma} does in base 3.
  \item \textbf{Structure of the table.} Explain the set $E$ of entry offsets and the
        non-uniformities in Table~\ref{tab:cr} directly.
  \item \textbf{Lengths.} The table gives the exact number of terms of every passage as an affine
        function of $b^k$. Use it to study the lengths (1.2) of~\cite{comma} for all starts.
  \item \textbf{OEIS.} The terms of \oeis{A399179}, and the uniqueness of the infinite path, could be
        contributed to the OEIS.
\end{itemize}
